%%%%%
%
%% Dans ce fichier, les �ventuels accents sont cod�s
%% en ISO-8859-15 (aka latin9)
%% La ligne suivante donne trois exemples 
%% e accent aigu : �, e accent grave : �, e-dans-l'o : �
\section{�chantillonnage et moyenne harmonique d'une variable sur une nouvelle base
de temps donn�e par un time\_bnds}

\itemstyle{Op\'erateur~:} \operfmt{MoyHarmEch\_tps\_bnds}
\index{MoyHarmEch\_tps\_bnds}

%%%%%OBJET%%%%
\itemstyle{Objet~:}
Cette fonction calcule, pour chaque instant $t$ d'une
base de temps $time_{out}$, 
une variable $X_{out}$, contenant la moyenne harmonique des
valeurs de  $X_{in}$ entre les bornes donn�es par le $time_{bnds}$ associ� � $time_{out}$.

Aucune supposition n'est faite sur la r�gularit� de
l'�chantillonnage des 2 variables. 
Il est cependant n�cessaire que $time_{in}$ (le temps de $X_{in}$) 
soit tri� par ordre croissant.


%%%%%ARGUMENTS%%%%
\bigskip
%% Dans ce fichier, les �ventuels accents sont cod�s
%% en ISO-8859-15 (aka latin9)
%% La ligne suivante donne trois exemples 
%% e accent aigu : �, e accent grave : �, e-dans-l'o : �
%%%%%ENTREES%%%%
\begin{tabarg}
\multicolumn{3}{l}{\titretab {Entr\'ees}}\\
\hline
$X_{in}$ & Vecteur, \texttt{float} & Variable d'origine\\
$toler_{mnq}$ & Coeff., \texttt{Int\_32} & Seuil de calcul de la moyenne (nb. �chant). A chaque instant de la base de temps d'arriv�e, 
la moyenne harmonique est calcul�e si le nombre d'�chantillons dans l'intervalle de temps est strictement sup�rieur � $toler_{mnq}$. Sinon, 
le r�sultat est "valeur manquante". Si $toler_{mnq}$ = 0, le calcul est fait d�s qu'il y a au moins une valeur valide dans l'�chantillon. \\
$time_{out}$ & Vecteur, \texttt{double} & Base de temps d'arriv�e (secondes)\\
\end{tabarg}

\medskip
%%%%%SORTIES%%%%
\begin{tabarg}
\multicolumn{3}{l}{\titretab {Sorties}}\\
\hline
$X_{out}$ & Vecteur, \texttt{float} & $X_{in}$ moyenn� et �chantillonn� sur
la base de $time_{out}$ (m�me unit� que $X_{in})$\\
\end{tabarg}



%%%%%FORMULE%%%%
\itemstyle{Formule}

%% Formule de l'operateur : MoyHarmEch_tps_bnds
%% Dans ce fichier, les �ventuels accents sont cod�s
%% en ISO-8859-15 (aka latin9)
%% La ligne suivante donne trois exemples 
%% e accent aigu : �, e accent grave : �, e-dans-l'o : �
Pour  i=$0$, Nbre $Time_{bnds}$ :
\[
X_{out}[i] = \frac{N_{i}}
{\displaystyle{\sum_{tb[i][0]}^{tb[i][1]}  \frac{1}{X_{in}[j]}} }
\]
 
Pour chaque $X_{in}[j]$ diff\'{e}rent de NaN avec $N_i$ le nombre de donn\'ees dans l'intervalle tb  pour lesquels $X_{in}[j]$ est diff\'erent de NaN.


%%%%COMMENTAIRE%%%%
\itemstyle{Commentaire~:}
\textit{tb} est un tableau d'intervalles de temps d�termin�s � partir du vecteur temps $time_{out}$.

%%%%%AUTEUR%%%%
\itemstyle{Auteur~:}
CNRM/GMEI/TRAMM

%%%%%REFERENCE%%%%
\itemstyle{R\'ef\'erences~:}
N�ant.
